\examplebreak
\subsubsection*{Ví dụ Adams 2: đủ 1 bước khởi tạo và một bước Adams}
Ví dụ tự dựng: $u^{\prime}=u$, $u(0)=1$, $h=0.1$. Gọi $R_4=265241/240000$ là hệ số tăng trưởng một bước RK4. Ở bước Adams dưới đây, $n=s-1$.
\begin{center}\begin{tabular}{rrll}\toprule $i$&$t_i$&$u_i$&Cách có giá trị\\\midrule
0&0&1&Điều kiện đầu\\
1&0.1&1.10517083333&RK4: $u_1 = R_4u_{0}$\\
\bottomrule\end{tabular}\end{center}
Ở bước Adams, xếp lịch sử từ mới nhất đến cũ nhất:
\begin{center}\begin{tabular}{rlll}\toprule $j$&$f_{n-j}$&$\beta_jf_{n-j}$&$a_{j+1}f_{n-j}$\\\midrule
0&1.10517083333&1.65775625&0.736780555556\\
1&1&-0.5&-0.0833333333333\\
\bottomrule\end{tabular}\end{center}
Dự báo: $u_{n+1}^P=u_n+0.1\sum\beta_jf_{n-j}=1.22094645833$.
Hiệu chỉnh giải phương trình tuyến tính $(1-0.1\,5/12)u_{n+1}=u_n+0.1\sum a_{j+1}f_{n-j}$.
Kết quả: $u_{n+1}=1.22140753623$ tại $t_{n+1}=0.2$. Đối chiếu app: $1.22140753623$.
Đây là phép giải đại số của phương trình hiệu chỉnh trong ví dụ tuyến tính, không phải nhật ký các lần gọi nội bộ của bộ giải phi tuyến.
\examplebreak
\subsubsection*{Ví dụ Adams 3: đủ 2 bước khởi tạo và một bước Adams}
Ví dụ tự dựng: $u^{\prime}=u$, $u(0)=1$, $h=0.1$. Gọi $R_4=265241/240000$ là hệ số tăng trưởng một bước RK4. Ở bước Adams dưới đây, $n=s-1$.
\begin{center}\begin{tabular}{rrll}\toprule $i$&$t_i$&$u_i$&Cách có giá trị\\\midrule
0&0&1&Điều kiện đầu\\
1&0.1&1.10517083333&RK4: $u_1 = R_4u_{0}$\\
2&0.2&1.22140257085&RK4: $u_2 = R_4u_{1}$\\
\bottomrule\end{tabular}\end{center}
Ở bước Adams, xếp lịch sử từ mới nhất đến cũ nhất:
\begin{center}\begin{tabular}{rlll}\toprule $j$&$f_{n-j}$&$\beta_jf_{n-j}$&$a_{j+1}f_{n-j}$\\\midrule
0&1.22140257085&2.34102159413&0.966943701923\\
1&1.10517083333&-1.47356111111&-0.230243923611\\
2&1&0.416666666667&0.0416666666667\\
\bottomrule\end{tabular}\end{center}
Dự báo: $u_{n+1}^P=u_n+0.1\sum\beta_jf_{n-j}=1.34981528582$.
Hiệu chỉnh giải phương trình tuyến tính $(1-0.1\,3/8)u_{n+1}=u_n+0.1\sum a_{j+1}f_{n-j}$.
Kết quả: $u_{n+1}=1.34985892504$ tại $t_{n+1}=0.3$. Đối chiếu app: $1.34985892504$.
Đây là phép giải đại số của phương trình hiệu chỉnh trong ví dụ tuyến tính, không phải nhật ký các lần gọi nội bộ của bộ giải phi tuyến.
\examplebreak
\subsubsection*{Ví dụ Adams 4: đủ 3 bước khởi tạo và một bước Adams}
Ví dụ tự dựng: $u^{\prime}=u$, $u(0)=1$, $h=0.1$. Gọi $R_4=265241/240000$ là hệ số tăng trưởng một bước RK4. Ở bước Adams dưới đây, $n=s-1$.
\begin{center}\begin{tabular}{rrll}\toprule $i$&$t_i$&$u_i$&Cách có giá trị\\\midrule
0&0&1&Điều kiện đầu\\
1&0.1&1.10517083333&RK4: $u_1 = R_4u_{0}$\\
2&0.2&1.22140257085&RK4: $u_2 = R_4u_{1}$\\
3&0.3&1.34985849706&RK4: $u_3 = R_4u_{2}$\\
\bottomrule\end{tabular}\end{center}
Ở bước Adams, xếp lịch sử từ mới nhất đến cũ nhất:
\begin{center}\begin{tabular}{rlll}\toprule $j$&$f_{n-j}$&$\beta_jf_{n-j}$&$a_{j+1}f_{n-j}$\\\midrule
0&1.34985849706&3.09342572243&1.21112304042\\
1&1.22140257085&-3.00261465334&-0.447847609312\\
2&1.10517083333&1.70380503472&0.162705706019\\
3&1&-0.375&-0.0263888888889\\
\bottomrule\end{tabular}\end{center}
Dự báo: $u_{n+1}^P=u_n+0.1\sum\beta_jf_{n-j}=1.49182010744$.
Hiệu chỉnh giải phương trình tuyến tính $(1-0.1\,251/720)u_{n+1}=u_n+0.1\sum a_{j+1}f_{n-j}$.
Kết quả: $u_{n+1}=1.49182437726$ tại $t_{n+1}=0.4$. Đối chiếu app: $1.49182437726$.
Đây là phép giải đại số của phương trình hiệu chỉnh trong ví dụ tuyến tính, không phải nhật ký các lần gọi nội bộ của bộ giải phi tuyến.
\examplebreak
\subsubsection*{Ví dụ Adams 5: đủ 4 bước khởi tạo và một bước Adams}
Ví dụ tự dựng: $u^{\prime}=u$, $u(0)=1$, $h=0.1$. Gọi $R_4=265241/240000$ là hệ số tăng trưởng một bước RK4. Ở bước Adams dưới đây, $n=s-1$.
\begin{center}\begin{tabular}{rrll}\toprule $i$&$t_i$&$u_i$&Cách có giá trị\\\midrule
0&0&1&Điều kiện đầu\\
1&0.1&1.10517083333&RK4: $u_1 = R_4u_{0}$\\
2&0.2&1.22140257085&RK4: $u_2 = R_4u_{1}$\\
3&0.3&1.34985849706&RK4: $u_3 = R_4u_{2}$\\
4&0.4&1.49182424008&RK4: $u_4 = R_4u_{3}$\\
\bottomrule\end{tabular}\end{center}
Ở bước Adams, xếp lịch sử từ mới nhất đến cũ nhất:
\begin{center}\begin{tabular}{rlll}\toprule $j$&$f_{n-j}$&$\beta_jf_{n-j}$&$a_{j+1}f_{n-j}$\\\midrule
0&1.49182424008&3.93883038944&1.47835638236\\
1&1.34985849706&-5.20070482063&-0.748046583789\\
2&1.22140257085&4.43776267409&0.408830582743\\
3&1.10517083333&-1.9555383912&-0.132773995949\\
4&1&0.348611111111&0.01875\\
\bottomrule\end{tabular}\end{center}
Dự báo: $u_{n+1}^P=u_n+0.1\sum\beta_jf_{n-j}=1.64872033636$.
Hiệu chỉnh giải phương trình tuyến tính $(1-0.1\,95/288)u_{n+1}=u_n+0.1\sum a_{j+1}f_{n-j}$.
Kết quả: $u_{n+1}=1.64872076496$ tại $t_{n+1}=0.5$. Đối chiếu app: $1.64872076496$.
Đây là phép giải đại số của phương trình hiệu chỉnh trong ví dụ tuyến tính, không phải nhật ký các lần gọi nội bộ của bộ giải phi tuyến.
